\documentclass[UTF8,11pt,a4paper]{ctexart}
\usepackage[margin=1.8cm]{geometry}
\usepackage{amsmath,amssymb,mathtools,booktabs,array,longtable}
\usepackage{enumitem}
\usepackage{needspace}
\usepackage{xcolor}
\usepackage{microtype}
\setlength{\parindent}{0pt}
\setlength{\parskip}{0.35em}
\linespread{1.12}
\newcommand{\Pp}{\mathrm{P}}
\newcommand{\E}{\mathrm{E}}
\newcommand{\D}{\mathrm{D}}
\newcommand{\Cov}{\mathrm{Cov}}
\newcommand{\Cnk}[2]{\mathrm{C}_{#1}^{#2}}
\newcommand{\Var}{\mathrm{Var}}
\newcommand{\Ans}[1]{\par\textbf{答案：}#1\par}
\newcommand{\Ana}[1]{\par\textbf{解析：}#1\par\medskip}
\newcommand{\Opt}[4]{\par A.\,#1\qquad B.\,#2\par C.\,#3\qquad D.\,#4\par}
\newcommand{\Q}[1]{\par\textbf{#1}\par}
\newcommand{\ruleline}{\par\noindent\rule{\linewidth}{0.35pt}\par}
\begin{document}

\begin{center}
{\LARGE\bfseries 概率论与数理统计 A\par 复习卷题目与解析}\vspace{0.4em}
\end{center}

\section*{复习卷（一）}
\subsection*{1. 单项选择题}

\Q{1. 设 $A,B$ 互为对立事件，且 $\Pp(A)>0,\Pp(B)>0$，则下列结论不正确的是（\quad）。}
\Opt{$\Pp(B)=1-\Pp(A)$}{$\Pp(A\cup B)=\Pp(A)+\Pp(B)$}{$\Pp(A\mid B)=0$}{$\Pp(\overline{A\cup B})=1$}
\Ans{D}
\Ana{对立事件满足 $A\cup B=\Omega$，故 $\Pp(\overline{A\cup B})=0$，不是 $1$。}

\Q{2. 已知随机变量 $X$ 服从泊松分布，且 $\D(X)=1$，则 $\Pp(X=1)=$（\quad）。}
\Opt{$2e$}{$0$}{$e$}{$e^{-1}$}
\Ans{D}
\Ana{泊松分布有 $\D(X)=\lambda=1$，所以 $\Pp(X=1)=e^{-1}\dfrac{1^1}{1!}=e^{-1}$。}

\Q{3. 设随机变量 $X$ 的分布函数为 $F(x)$，下列不一定正确的是（\quad）。}
\Opt{$F(x)$ 必为连续函数}{$\lim\limits_{x\to-\infty}F(x)=0$}{$\Pp(0<X\le1)=F(1)-F(0)$}{$F(0)\le F(1)$}
\Ans{A}
\Ana{分布函数只要求右连续、单调不减，不一定处处连续。}

\Q{4. 若随机变量 $X$ 满足 $\D(2X+1)=10$，则 $\D(X)=$（\quad）。}
\Opt{$2.25$}{$2.5$}{$4.5$}{$5$}
\Ans{B}
\Ana{$\D(2X+1)=2^2\D(X)=10$，故 $\D(X)=\dfrac{10}{4}=2.5$。}

\Q{5. 若随机变量 $X,Y$ 相互独立，且 $X\sim N(1,4^2)$，$Y\sim N(2,1^2)$，则 $\E(X-2Y+3)=$（\quad）。}
\Opt{$12$}{$9$}{$0$}{$-3$}
\Ans{C}
\Ana{$\E(X-2Y+3)=1-2\times2+3=0$。}

\Q{6. 设 $X_1,X_2,X_3,X_4$ 是来自总体 $X\sim N(\mu,\sigma^2)$ 的样本，其中 $\mu$ 已知、$\sigma$ 未知，则下列样本函数中能作为统计量的是（\quad）。}
\Opt{$\max X_i$}{$\displaystyle\sum_{i=1}^4\frac{X_i-\mu}{\sigma}$}{$\displaystyle\frac1{\sigma^2}\sum_{i=1}^4X_i^2$}{$\mu X_1+\sigma X_2$}
\Ans{A}
\Ana{统计量不得包含未知参数。B、C、D 都含未知参数 $\sigma$；A 仅由样本构成。}

\subsection*{2. 填空题}
\Q{1. 设随机事件满足 $A\subset B$，$\Pp(A)=0.3$，$\Pp(B)=0.5$，则 $\Pp(B-A)=$\underline{\hspace{2cm}}。}
\Ans{$0.2$}
\Ana{$B=A\cup(B-A)$ 且两部分互不相容，故 $\Pp(B-A)=0.5-0.3=0.2$。}

\Q{2. 设 $100$ 件产品中有 $5$ 件次品，现不放回地连续随机抽取两次，每次抽取 $1$ 件，则第二次取到次品的概率为\underline{\hspace{2cm}}。}
\Ans{$\dfrac5{100}$}
\Ana{任意一个抽取位置取到次品的边缘概率均为总体次品比例 $\dfrac5{100}$。}

\Q{3. 设 $X$ 为连续型随机变量，则 $\Pp(X=1)=$\underline{\hspace{2cm}}。}
\Ans{$0$}
\Ana{连续型随机变量在单点处的概率为 $0$。}

\Q{4. 从区间 $[0,1]$ 上随机取出一个数，记为 $X$，则 $\Pp(4X-1\le0)=$\underline{\hspace{2cm}}。}
\Ans{$\dfrac14$}
\Ana{$4X-1\le0\Leftrightarrow X\le\dfrac14$，均匀分布下概率为区间长度 $\dfrac14$。}

\Q{5. 设随机变量 $X,Y$ 独立同分布，且 $\Pp(X=1)=\dfrac14$，$\Pp(X=2)=\dfrac34$，则 $\Pp(X+Y=3)=$\underline{\hspace{2cm}}。}
\Ans{$\dfrac38$}
\Ana{$\Pp(X+Y=3)=2\Pp(X=1)\Pp(Y=2)=2\times\dfrac14\times\dfrac34=\dfrac38$。}

\Q{6. 若随机变量 $X\sim B\left(2,\dfrac12\right)$，则 $\E(X^2)=$\underline{\hspace{2cm}}。}
\Ans{$\dfrac32$}
\Ana{$\E(X)=1$，$\D(X)=\dfrac12$，所以 $\E(X^2)=\D(X)+[\E(X)]^2=\dfrac12+1=\dfrac32$。}

\Q{7. 设随机变量 $(X,Y)$ 的联合概率密度为 $f(x,y)$，则 $\displaystyle\int_{-\infty}^{+\infty}\int_{-\infty}^{+\infty}f(x,y)\,dx\,dy=$\underline{\hspace{2cm}}。}
\Ans{$1$}
\Ana{联合概率密度在全平面上的积分恒为 $1$。}

\Q{8. 设随机变量 $X$ 的分布律如表，其中 $\theta$ 是未知参数，$\overline X$ 是样本均值，则 $\theta$ 的矩估计量为\underline{\hspace{2cm}}。}
\[
\begin{array}{c|ccc}X&0&1&2\\\hline \Pp&1-3\theta&\theta&2\theta\end{array}
\]
\Ans{$\hat\theta=\dfrac{\overline X}{5}$}
\Ana{$\E(X)=0(1-3\theta)+\theta+2(2\theta)=5\theta$。令 $\overline X=5\theta$ 得 $\hat\theta=\overline X/5$。}

\Needspace{16\baselineskip}\subsection*{3. 计算题}
\Q{1. 箱中有 $8$ 个白球和 $2$ 个黑球，从中随机抽取 $2$ 个球。每取出 $1$ 个白球得 $1$ 分，每取出 $1$ 个黑球扣 $1$ 分。用 $X$ 表示得到的分数，求：（1）$X$ 的分布律；（2）$X$ 的数学期望。}
\Ana{
设抽到白球的个数为 $W$。因为共抽取 $2$ 个球，所以抽到黑球的个数为 $2-W$，得分为
\[
X=W-(2-W)=2W-2.
\]
因此，$W=0,1,2$ 时，$X$ 分别取 $-2,0,2$。所有等可能的抽法总数为
\[
\Cnk{10}{2}=45.
\]
\textbf{（1）求分布律：}
\[
\begin{aligned}
\Pp(X=-2)
&=\Pp(W=0)
 =\frac{\Cnk{2}{2}\Cnk{8}{0}}{\Cnk{10}{2}}
 =\frac{1}{45},\\
\Pp(X=0)
&=\Pp(W=1)
 =\frac{\Cnk{8}{1}\Cnk{2}{1}}{\Cnk{10}{2}}
 =\frac{16}{45},\\
\Pp(X=2)
&=\Pp(W=2)
 =\frac{\Cnk{8}{2}\Cnk{2}{0}}{\Cnk{10}{2}}
 =\frac{28}{45}.
\end{aligned}
\]
故 $X$ 的分布律为
\[
\begin{array}{c|ccc}
X&-2&0&2\\\hline
\Pp&\dfrac1{45}&\dfrac{16}{45}&\dfrac{28}{45}
\end{array}
\]
\textbf{（2）求数学期望：}
\[
\begin{aligned}
\E(X)
&=(-2)\times\frac1{45}+0\times\frac{16}{45}+2\times\frac{28}{45}\\
&=\frac{-2+56}{45}=\frac{54}{45}=\frac65.
\end{aligned}
\]}

\Q{2. 设二维随机变量 $(X,Y)$ 的联合分布律如表，且 $\Pp(X=1\mid Y=0)=\dfrac35$，求：（1）$a,b$ 的值；（2）$X,Y$ 的边缘分布律，并判断 $X,Y$ 是否独立。}
\[
\begin{array}{c|ccc}
 &Y=0&Y=1&Y=2\\\hline
X=0&\dfrac2{25}&a&\dfrac1{25}\\
X=1&b&\dfrac3{25}&\dfrac2{25}
\end{array}
\]
\Ana{
\textbf{（1）求 $a,b$：}

由条件概率定义，先求
\[
\Pp(Y=0)=\Pp(X=0,Y=0)+\Pp(X=1,Y=0)=\frac2{25}+b.
\]
于是
\[
\begin{aligned}
\Pp(X=1\mid Y=0)
&=\frac{\Pp(X=1,Y=0)}{\Pp(Y=0)}\\
&=\frac{b}{\frac2{25}+b}=\frac35.
\end{aligned}
\]
解得
\[
5b=3\left(\frac2{25}+b\right)
\Rightarrow 2b=\frac6{25}
\Rightarrow b=\frac3{25}.
\]
再由联合分布律中所有概率之和为 $1$：
\[
\frac2{25}+a+\frac1{25}+\frac3{25}+\frac3{25}+\frac2{25}=1,
\]
即
\[
a+\frac{11}{25}=1\Rightarrow a=\frac{14}{25}.
\]

\textbf{（2）求边缘分布律并判断独立性：}
\[
\begin{aligned}
\Pp(X=0)&=\frac2{25}+\frac{14}{25}+\frac1{25}=\frac{17}{25},\\
\Pp(X=1)&=\frac3{25}+\frac3{25}+\frac2{25}=\frac8{25}.
\end{aligned}
\]
\[
\begin{aligned}
\Pp(Y=0)&=\frac2{25}+\frac3{25}=\frac5{25},\\
\Pp(Y=1)&=\frac{14}{25}+\frac3{25}=\frac{17}{25},\\
\Pp(Y=2)&=\frac1{25}+\frac2{25}=\frac3{25}.
\end{aligned}
\]
所以边缘分布律为
\[
\begin{array}{c|cc}
X&0&1\\\hline
\Pp&\dfrac{17}{25}&\dfrac8{25}
\end{array}
\qquad
\begin{array}{c|ccc}
Y&0&1&2\\\hline
\Pp&\dfrac5{25}&\dfrac{17}{25}&\dfrac3{25}
\end{array}
\]
检验一组取值：
\[
\Pp(X=0,Y=0)=\frac2{25},
\]
而
\[
\Pp(X=0)\Pp(Y=0)=\frac{17}{25}\cdot\frac5{25}\ne\frac2{25}.
\]
故 $X,Y$ 不独立。
}

\subsection*{4. 解答题}
\Q{1. 已知区域 $G=\{(x,y)\mid0\le x\le1,\ 0\le y\le1\}$，设二维随机变量 $(X,Y)$ 在区域 $G$ 上服从均匀分布。求：（1）$(X,Y)$ 的联合概率密度；（2）$\Pp\left(X<\dfrac12\right)$；（3）$\Pp\left(Y<\dfrac23\mid X<\dfrac12\right)$。}
\Ana{区域面积 $S_G=1$，故
\[
f(x,y)=\begin{cases}1,&(x,y)\in G,\\0,&\text{其他}。\end{cases}
\]
\[
\Pp\left(X<\frac12\right)=\frac{\frac12\times1}{1}=\frac12.
\]
\[
\Pp\left(Y<\frac23\mid X<\frac12\right)=\frac{\frac12\times\frac23}{\frac12}=\frac23.
\]}

\Q{2. 设总体 $X$ 的概率密度为
\[
f(x;\theta)=\begin{cases}(\theta-2)x^{\theta-3},&0<x<1,\\0,&\text{其他},\end{cases}\qquad \theta>2,
\]
其中 $\theta$ 未知，$X_1,\ldots,X_n$ 是总体 $X$ 的样本。用极大似然估计法估计 $\theta$。}
\Ana{
\[
L(\theta)=(\theta-2)^n\left(\prod_{i=1}^n x_i\right)^{\theta-3},
\]
\[
\ell(\theta)=n\ln(\theta-2)+(\theta-3)\ln\left(\prod_{i=1}^n x_i\right).
\]
\[
\ell'(\theta)=\frac{n}{\theta-2}+\ln\left(\prod_{i=1}^n x_i\right)=0.
\]
\[
\hat\theta=2-\frac{n}{\ln\left(\prod_{i=1}^nX_i\right)}.
\]}

\subsection*{5. 应用题}
\Q{1. 某电商有甲、乙、丙三家合作快递，选择三家快递发货的概率分别为 $0.5,0.3,0.2$；甲、乙、丙的准时送达率分别为 $0.9,0.8,0.95$。求：（1）该店发货能准时送达的概率；（2）已知某次发货准时送达，求其选择甲快递的概率。}
\Ana{令 $A_1,A_2,A_3$ 分别表示选择甲、乙、丙，$B$ 表示准时送达。
\[
\Pp(B)=0.5\times0.9+0.3\times0.8+0.2\times0.95=0.88.
\]
\[
\Pp(A_1\mid B)=\frac{0.5\times0.9}{0.88}=\frac{45}{88}.
\]}

\Q{2. 某品牌洗衣机的使用寿命 $X$（单位：年）服从正态分布 $N(10,0.1^2)$。求：（1）平均使用寿命；（2）使用寿命超过 $10.1$ 年的概率；（3）使用寿命在 $9.8$ 年至 $10.2$ 年的概率。已知 $\Phi(1)=0.8413$，$\Phi(2)=0.9772$。}
\Ana{
\[
\E(X)=10.
\]
\[
\Pp(X>10.1)=1-\Phi\left(\frac{10.1-10}{0.1}\right)=1-\Phi(1)=0.1587.
\]
\[
\Pp(9.8<X<10.2)=\Phi(2)-\Phi(-2)=2\Phi(2)-1=0.9544.
\]}

\subsection*{6. 证明题}
\Q{已知随机变量 $X$ 服从参数为 $\lambda$ 的指数分布，且满足 $\Pp(X>1)=e^{-2}$。（1）写出 $X$ 的分布函数，并求 $\lambda$；（2）证明 $\dfrac{\E(X)}{\D(X)}=2$。}
\Ana{指数分布有 $\Pp(X>1)=e^{-\lambda}$，故 $e^{-\lambda}=e^{-2}$，即 $\lambda=2$。
\[
F(x)=\begin{cases}1-e^{-2x},&x>0,\\0,&x\le0.\end{cases}
\]
\[
\E(X)=\frac12,\qquad \D(X)=\frac14,\qquad \frac{\E(X)}{\D(X)}=\frac{1/2}{1/4}=2.
\]}

\newpage
\section*{复习卷（二）}
\subsection*{1. 单项选择题}

\Q{1. 设 $A,B$ 均为非零概率事件，且 $A,B$ 互不相容，则下列结论中正确的是（\quad）。}
\Opt{$\Pp(A)=1-\Pp(B)$}{$\Pp(A\cap B)=\Pp(A)\Pp(B)$}{$\overline A$ 与 $B$ 亦互不相容}{$\Pp(A\cup B)=\Pp(A)+\Pp(B)$}
\Ans{D}
\Ana{互不相容事件满足 $A\cap B=\varnothing$，因此并事件概率等于概率之和。}

\Q{2. 设随机变量 $X\sim B(3,0.3)$，则 $\Pp(X>2)=$（\quad）。}
\Opt{$0.027$}{$0.09$}{$0.3$}{$1$}
\Ans{A}
\Ana{$\Pp(X>2)=\Pp(X=3)=0.3^3=0.027$。}

\Q{3. 若随机变量 $X,Y$ 相互独立，且 $X\sim N(-1,3^2)$，$Y\sim N(1,2^2)$，则 $\D(X-Y)=$（\quad）。}
\Opt{$5$}{$13$}{$0$}{$1$}
\Ans{B}
\Ana{独立时 $\D(X-Y)=\D(X)+\D(Y)=3^2+2^2=13$。}

\Q{4. 设 $X,Y$ 为任意随机变量，则下列各式一定成立的是（\quad）。}
\Opt{$\E(XY)=\E(X)\E(Y)$}{$\E(X+Y)=\E(X)+\E(Y)$}{$\D(XY)=\D(X)\D(Y)$}{$\D(X-Y)=\D(X)+\D(Y)$}
\Ans{B}
\Ana{期望的线性性质总成立；其余各式通常需要独立等附加条件，且 C 即使独立也不一般成立。}

\Q{5. 设随机变量 $X$ 的数学期望 $\E(X)=-2$，方差 $\D(X)=0$，则 $\E(X^2)=$（\quad）。}
\Opt{$0$}{$2$}{$4$}{$-2$}
\Ans{C}
\Ana{$\E(X^2)=\D(X)+[\E(X)]^2=0+(-2)^2=4$。}

\Q{6. 设 $X_1,\ldots,X_n$ 是来自总体 $X\sim N(\mu,\sigma^2)$ 的样本，$\overline X$ 为样本均值，则下列关于样本方差 $S^2$ 的说法正确的是（\quad）。}
\Opt{$\displaystyle S^2=\frac1n\sum_{i=1}^n(X_i-\sigma)^2$}{$\displaystyle S^2=\frac1n\sum_{i=1}^n(X_i-\overline X)^2$}{$\displaystyle S^2=\frac1{n-1}\sum_{i=1}^n(X_i-\sigma)^2$}{$\displaystyle S^2=\frac1{n-1}\sum_{i=1}^n(X_i-\overline X)^2$}
\Ans{D}
\Ana{通常定义的无偏样本方差为 $S^2=\dfrac1{n-1}\sum_{i=1}^n(X_i-\overline X)^2$。}

\subsection*{2. 填空题}
\Q{1. 已知 $5$ 件产品中有 $2$ 件一等品、$3$ 件二等品，从中任取 $3$ 件，则恰取出 $2$ 件一等品的概率为\underline{\hspace{2cm}}。}
\Ans{$\dfrac3{10}$}
\Ana{$\dfrac{\Cnk{2}{2}\Cnk{3}{1}}{\Cnk{5}{3}}=\dfrac3{10}$。}

\Q{2. 设随机变量 $X\sim B\left(n,\dfrac13\right)$，且 $\D(X)=8$，则 $n=$\underline{\hspace{2cm}}。}
\Ans{$36$}
\Ana{$\D(X)=n\cdot\dfrac13\cdot\dfrac23=\dfrac{2n}{9}=8$，故 $n=36$。}

\Q{3. 设随机事件 $A,B$ 相互独立，$\Pp(A)=\Pp(B)=\dfrac23$，则 $\Pp(\overline A\,\overline B)=$\underline{\hspace{2cm}}。}
\Ans{$\dfrac19$}
\Ana{$\Pp(\overline A\,\overline B)=\Pp(\overline A)\Pp(\overline B)=\dfrac13\cdot\dfrac13=\dfrac19$。}

\Q{4. 设 $F(x)$ 为随机变量 $X$ 的分布函数，且 $\Pp(X>1)=0.15$，则 $F(1)=$\underline{\hspace{2cm}}。}
\Ans{$0.85$}
\Ana{$F(1)=\Pp(X\le1)=1-\Pp(X>1)=0.85$。}

\Q{5. 设 $X,Y$ 为随机变量，且 $\D(X+Y)=7$，$\D(X)=4$，$\D(Y)=1$，则 $\Cov(X,Y)=$\underline{\hspace{2cm}}。}
\Ans{$1$}
\Ana{$7=4+1+2\Cov(X,Y)$，故 $\Cov(X,Y)=1$。}

\Q{6. 设随机变量 $X,Y$ 相互独立，且 $\Pp(X\le1)=\dfrac12$，$\Pp(Y\le1)=\dfrac13$，则 $\Pp(X\le1,Y\le1)=$\underline{\hspace{2cm}}。}
\Ans{$\dfrac16$}
\Ana{独立性给出 $\dfrac12\cdot\dfrac13=\dfrac16$。}

\Q{7. 设随机变量 $X$ 服从区间 $[1,5]$ 上的均匀分布，则 $\Pp(2<X<3)=$\underline{\hspace{2cm}}。}
\Ans{$\dfrac14$}
\Ana{概率为区间长度比：$\dfrac{3-2}{5-1}=\dfrac14$。}

\Q{8. 已知某电子仪器的使用寿命 $X\sim N(\mu,2)$，现随机抽取 $4$ 只，测得寿命分别为 $0.85,1.1,0.9,1.15$，则 $\mu$ 的矩估计值为\underline{\hspace{2cm}}。}
\Ans{$1$}
\Ana{正态总体均值的矩估计为样本均值：$\hat\mu=\overline X=\dfrac{0.85+1.1+0.9+1.15}{4}=1$。}

\Needspace{16\baselineskip}\subsection*{3. 计算题}
\Q{1. 已知正方形的边长 $X$（m）是一个随机变量，其分布律如下：
\[
\begin{array}{c|ccc}X&\frac14&\frac12&1\\\hline \Pp&0.5&0.25&0.25\end{array}
\]
求：（1）正方形周长 $Y$ 的分布律；（2）$\E(Y^2)$。}
\Ana{
\textbf{（1）求 $Y$ 的分布律：}

正方形周长为
\[
Y=4X.
\]
逐个代入 $X$ 的可能取值：
\[
\begin{array}{c|ccc}
X&\dfrac14&\dfrac12&1\\\hline
Y=4X&1&2&4
\end{array}
\]
三个 $X$ 的取值映射到三个不同的 $Y$ 值，因此概率不需要合并：
\[
\begin{aligned}
\Pp(Y=1)&=\Pp\left(X=\frac14\right)=0.5,\\
\Pp(Y=2)&=\Pp\left(X=\frac12\right)=0.25,\\
\Pp(Y=4)&=\Pp(X=1)=0.25.
\end{aligned}
\]
故
\[
\begin{array}{c|ccc}
Y&1&2&4\\\hline
\Pp&0.5&0.25&0.25
\end{array}
\]
\textbf{（2）求 $\E(Y^2)$：}
\[
\begin{aligned}
\E(Y^2)
&=1^2\times0.5+2^2\times0.25+4^2\times0.25\\
&=0.5+1+4=5.5.
\end{aligned}
\]
}

\Q{2. 设二维随机变量 $(X,Y)$ 的联合分布律如表，且 $\Pp(X=0)=0.5$。求：（1）$a,b$ 的值；（2）$X,Y$ 的协方差，并判断 $X,Y$ 是否相关。}
\[
\begin{array}{c|ccc}&Y=-1&Y=0&Y=1\\\hline
X=0&a&0.1&0.2\\
X=1&0.1&b&0.2
\end{array}
\]
\Ana{
\textbf{（1）求 $a,b$：}

由边缘概率 $\Pp(X=0)=0.5$，有
\[
a+0.1+0.2=0.5\Rightarrow a=0.2.
\]
再由全部联合概率之和为 $1$：
\[
0.2+0.1+0.2+0.1+b+0.2=1
\Rightarrow b=0.2.
\]

\textbf{（2）求协方差：}

先求边缘分布：
\[
\begin{array}{c|cc}
X&0&1\\\hline
\Pp&0.2+0.1+0.2=0.5&0.1+0.2+0.2=0.5
\end{array}
\]
\[
\begin{array}{c|ccc}
Y&-1&0&1\\\hline
\Pp&0.2+0.1=0.3&0.1+0.2=0.3&0.2+0.2=0.4
\end{array}
\]
于是
\[
\begin{aligned}
\E(X)&=0\times0.5+1\times0.5=0.5,\\
\E(Y)&=(-1)\times0.3+0\times0.3+1\times0.4=0.1.
\end{aligned}
\]
计算 $\E(XY)$ 时，$X=0$ 的各项均为 $0$，故
\[
\begin{aligned}
\E(XY)
&=1\times(-1)\times0.1+1\times0\times0.2+1\times1\times0.2\\
&=-0.1+0+0.2=0.1.
\end{aligned}
\]
因此
\[
\Cov(X,Y)=\E(XY)-\E(X)\E(Y)=0.1-0.5\times0.1=0.05\ne0.
\]
故 $X,Y$ 相关。
}

\subsection*{4. 解答题}
\Q{1. 设二维随机变量 $(X,Y)$ 在平面区域 $D$ 上服从均匀分布，其中 $D$ 是以原点为圆心、半径为 $1$ 的圆域。求：（1）$(X,Y)$ 的联合概率密度函数；（2）$\Pp(Y>0)$；（3）$\Pp(X>0\mid Y>0)$。}
\Ana{区域面积 $S_D=\pi$，故
\[
f(x,y)=\begin{cases}\dfrac1\pi,&x^2+y^2\le1,\\0,&\text{其他}。\end{cases}
\]
上半圆面积为 $\pi/2$，所以 $\Pp(Y>0)=\dfrac12$。
\[
\Pp(X>0\mid Y>0)=\frac{\pi/4}{\pi/2}=\frac12.
\]}

\Q{2. 设总体 $X$ 的概率密度为
\[
f(x;\theta)=\begin{cases}\dfrac{\theta}{x^{\theta+1}},&x\ge1,\\0,&x<1,\end{cases}\qquad\theta>1,
\]
其中 $\theta$ 未知，$X_1,\ldots,X_n$ 是总体 $X$ 的样本。用极大似然估计法估计 $\theta$。}
\Ana{
\[
L(\theta)=\theta^n\left(\prod_{i=1}^nX_i\right)^{-(\theta+1)},
\]
\[
\ell(\theta)=n\ln\theta-(\theta+1)\ln\left(\prod_{i=1}^nX_i\right).
\]
\[
\ell'(\theta)=\frac n\theta-\ln\left(\prod_{i=1}^nX_i\right)=0.
\]
\[
\hat\theta=\frac{n}{\ln\left(\prod_{i=1}^nX_i\right)}.
\]}

\subsection*{5. 应用题}
\Q{1. 某地区成年居民中偏胖者占 $10\%$，不胖不瘦者占 $80\%$，偏瘦者占 $10\%$；三类人患高血压的概率分别为 $20\%,10\%,5\%$。求：（1）该地区成年居民患高血压的概率；（2）已知某成年居民患有高血压，求其为偏胖者的概率。}
\Ana{设 $A_1,A_2,A_3$ 分别表示偏胖、不胖不瘦、偏瘦，$B$ 表示患高血压。
\[
\Pp(B)=0.1\times0.2+0.8\times0.1+0.1\times0.05=0.105.
\]
\[
\Pp(A_1\mid B)=\frac{0.1\times0.2}{0.105}=\frac4{21}.
\]}

\Q{2. 某高校男生身高 $X$（单位：m）服从正态分布 $N(1.7,0.08^2)$。求：（1）平均身高；（2）身高超过 $1.86$ m 的概率；（3）身高在 $1.62$ m 至 $1.78$ m 的概率。已知 $\Phi(1)=0.8413$，$\Phi(2)=0.9772$。}
\Ana{
\[
\E(X)=1.7.
\]
\[
\Pp(X>1.86)=1-\Phi\left(\frac{1.86-1.7}{0.08}\right)=1-\Phi(2)=0.0228.
\]
\[
\Pp(1.62\le X\le1.78)=\Phi(1)-\Phi(-1)=2\Phi(1)-1=0.6826.
\]}

\subsection*{6. 证明题}
\Q{已知随机变量 $X$ 服从参数为 $\lambda$ 的泊松分布，且满足 $\Pp(X=1)=\Pp(X=2)$。（1）证明 $\lambda=2$；（2）求 $\dfrac{\D(X+2)}{\E(X)}$。}
\Ana{
\[
\Pp(X=1)=e^{-\lambda}\lambda,\qquad \Pp(X=2)=e^{-\lambda}\frac{\lambda^2}{2}.
\]
由两者相等且 $\lambda>0$，得 $\lambda=2$。泊松分布满足 $\E(X)=\D(X)=\lambda=2$，又 $\D(X+2)=\D(X)=2$，所以
\[
\frac{\D(X+2)}{\E(X)}=\frac22=1.
\]}

\newpage
\section*{复习卷（三）}
\subsection*{1. 单项选择题}

\Q{1. 设 $A,B$ 互为对立事件，则下列结论中正确的是（\quad）。}
\Opt{$\Pp(A\cap B)=\Pp(A)\Pp(B)$}{$\Pp(A)=1-\Pp(B)$}{$\Pp(A\cap B)=\varnothing$}{$\Pp(A\cup B)=\Omega$}
\Ans{B}
\Ana{对立事件满足 $\Pp(A)=1-\Pp(B)$。C、D 将概率与集合相等同，表达不成立；A 一般不成立。}

\Q{2. 设随机变量 $X\sim U(2,4)$，则 $\Pp(2<X<3)=$（\quad）。}
\Opt{$\Pp(1.5<X<2.5)$}{$\Pp(2.5<X<3.5)$}{$\Pp(3.5<X<4.5)$}{$\Pp(4.5<X<5.5)$}
\Ans{B}
\Ana{$\Pp(2<X<3)=\dfrac12$。B 中与支持区间 $[2,4]$ 的交为 $(2.5,3.5)$，长度同为 $1$，概率也是 $\dfrac12$。}

\Q{3. 若随机变量 $X,Y$ 相互独立，且 $X\sim N(3,4)$，$Y\sim N(2,9)$，则 $Z=3X-Y\sim$（\quad）。}
\Opt{$N(7,21)$}{$N(7,27)$}{$N(7,45)$}{$N(11,45)$}
\Ans{C}
\Ana{$\E(Z)=3\times3-2=7$，$\D(Z)=3^2\times4+9=45$，故 $Z\sim N(7,45)$。}

\Q{4. 设 $(X,Y)$ 为二维随机变量，则下列各式一定成立的是（\quad）。}
\Opt{$\D(X-Y)=\D(X)+\D(Y)$}{$\Cov(5X,5Y)=5\Cov(X,Y)$}{$\E(XY)=\E(X)\E(Y)$}{$\E(X+Y)=\E(X)+\E(Y)$}
\Ans{D}
\Ana{期望可加性总成立；$\Cov(5X,5Y)=25\Cov(X,Y)$。}

\Q{5. 设随机变量 $X\sim B\left(10,\dfrac13\right)$，则 $\dfrac{\D(X)}{\E(X)}=$（\quad）。}
\Opt{$\dfrac23$}{$\dfrac13$}{$1$}{$\dfrac{10}{3}$}
\Ans{A}
\Ana{$\E(X)=\dfrac{10}{3}$，$\D(X)=10\cdot\dfrac13\cdot\dfrac23=\dfrac{20}{9}$，比值为 $\dfrac23$。}

\Q{6. 已知总体 $X$ 的分布中含有参数 $\mu,\sigma$，其中 $\mu$ 已知、$\sigma$ 未知，$X_1,\ldots,X_n$ 为总体 $X$ 的样本，则下列不是统计量的是（\quad）。}
\Opt{$\displaystyle\sum_{i=1}^n(X_i-\overline X)^2$}{$4\mu(X_1^2+X_2^2+X_3^2)$}{$\displaystyle\frac1{3\sigma^2}\sum_{i=1}^nX_i^2$}{$\displaystyle\frac{X_1^4+X_2^4+X_3^4}{6}$}
\Ans{C}
\Ana{C 含未知参数 $\sigma$，因此不是统计量；B 中 $\mu$ 已知。}

\subsection*{2. 填空题}
\Q{1. 已知袋中有 $5$ 个黑球、$3$ 个白球，从中任取 $4$ 个球，则恰有 $3$ 个白球的概率为\underline{\hspace{2cm}}。}
\Ans{$\dfrac1{14}$}
\Ana{$\dfrac{\Cnk{3}{3}\Cnk{5}{1}}{\Cnk{8}{4}}=\dfrac5{70}=\dfrac1{14}$。}

\Q{2. 设随机变量 $X\sim E(2)$，则 $\E(2X+10)=$\underline{\hspace{2cm}}。}
\Ans{$11$}
\Ana{$\E(X)=\dfrac12$，故 $\E(2X+10)=2\times\dfrac12+10=11$。}

\Q{3. 设 $\Pp(A)=\dfrac12$，$\Pp(B\mid A)=\dfrac12$，且 $\Pp(A\cup B)=\dfrac7{12}$，则 $\Pp(B)=$\underline{\hspace{2cm}}。}
\Ans{$\dfrac13$}
\Ana{$\Pp(A\cap B)=\Pp(A)\Pp(B\mid A)=\dfrac14$；$\dfrac7{12}=\dfrac12+\Pp(B)-\dfrac14$，故 $\Pp(B)=\dfrac13$。}

\Q{4. 设 $f(x)$ 为随机变量 $X$ 的概率密度，则 $\displaystyle\int_{-\infty}^{+\infty}f(x)\,dx=$\underline{\hspace{2cm}}。}
\Ans{$1$}
\Ana{概率密度在全实数轴上的积分为 $1$。}

\Q{5. 设 $X,Y$ 为随机变量，且 $\Cov(X,Y)=2.5$，$\D(X)=\D(Y)=5$，则相关系数 $\rho_{XY}=$\underline{\hspace{2cm}}。}
\Ans{$\dfrac12$}
\Ana{$\rho_{XY}=\dfrac{2.5}{\sqrt{5\times5}}=\dfrac12$。}

\Q{6. 设随机变量 $X,Y$ 独立同分布，且 $\Pp(X=1)=\dfrac12$，$\Pp(X=2)=\dfrac12$，则 $\Pp(XY=4)=$\underline{\hspace{2cm}}。}
\Ans{$\dfrac14$}
\Ana{$XY=4$ 只能由 $X=2,Y=2$ 产生，概率为 $\dfrac12\times\dfrac12=\dfrac14$。}

\Q{9. 设随机变量 $X\sim N(10,\sigma^2)$，已知 $\Pp(10<X\le20)=0.3$，则 $\Pp(0<X\le10)=$\underline{\hspace{2cm}}。}
\Ans{$0.3$}
\Ana{正态分布关于均值 $10$ 对称，区间 $(0,10]$ 与 $(10,20]$ 关于均值对称，概率相等。}

\Q{10. 已知总体 $X\sim P(\lambda)$，样本均值 $\overline x=6$，则 $\lambda$ 的矩估计值为\underline{\hspace{2cm}}。}
\Ans{$6$}
\Ana{泊松分布有 $\E(X)=\lambda$，令 $\overline X=\lambda$，得 $\hat\lambda=6$。}

\Needspace{16\baselineskip}\subsection*{3. 计算题}
\Q{1. 已知离散型随机变量 $X$ 的分布律如下：
\[
\begin{array}{c|cccc}X&-2&-1&0&1\\\hline \Pp&0.3&0.2&0.4&c\end{array}
\]
求：（1）常数 $c$ 的值；（2）$Y=X^2+1$ 的分布律；（3）$\Pp(Y<3)$。}
\Ana{
\textbf{（1）求 $c$：}

由分布律中全部概率之和等于 $1$：
\[
0.3+0.2+0.4+c=1,
\]
所以
\[
c=0.1.
\]

\textbf{（2）求 $Y=X^2+1$ 的分布律：}

分别代入 $X$ 的取值：
\[
\begin{array}{c|cccc}
X&-2&-1&0&1\\\hline
Y=X^2+1&5&2&1&2
\end{array}
\]
因此
\[
\begin{aligned}
\Pp(Y=1)&=\Pp(X=0)=0.4,\\
\Pp(Y=2)&=\Pp(X=-1)+\Pp(X=1)=0.2+0.1=0.3,\\
\Pp(Y=5)&=\Pp(X=-2)=0.3.
\end{aligned}
\]
故
\[
\begin{array}{c|ccc}
Y&1&2&5\\\hline
\Pp&0.4&0.3&0.3
\end{array}
\]

\textbf{（3）求 $\Pp(Y<3)$：}

满足 $Y<3$ 的取值为 $Y=1,2$，所以
\[
\Pp(Y<3)=\Pp(Y=1)+\Pp(Y=2)=0.4+0.3=0.7.
\]
}

\Q{2. 设二维离散型随机变量 $(X,Y)$ 满足 $\Pp(XY=0)=1$，且 $X,Y$ 的分布律如下：
\[
\begin{array}{c|ccc}X&-1&0&1\\\hline \Pp&\frac14&\frac12&\frac14\end{array}
\qquad
\begin{array}{c|cc}Y&0&1\\\hline \Pp&\frac12&\frac12\end{array}
\]
求：（1）$\Cov(X,Y)$；（2）$X,Y$ 的联合分布律。}
\Ana{
\textbf{（1）构造联合分布律：}

由 $\Pp(XY=0)=1$ 可知，当 $xy\ne0$ 时，对应的联合概率均为 $0$。特别地，$Y=1$ 时只能有 $X=0$，所以
\[
\Pp(X=-1,Y=1)=0,
\qquad
\Pp(X=1,Y=1)=0.
\]
又因为 $\Pp(Y=1)=\dfrac12$，故
\[
\Pp(X=0,Y=1)=\frac12.
\]
由 $X$ 的边缘分布，
\[
\Pp(X=-1,Y=0)=\Pp(X=-1)=\frac14,
\]
\[
\Pp(X=1,Y=0)=\Pp(X=1)=\frac14.
\]
最后，
\[
\Pp(X=0,Y=0)=\Pp(X=0)-\Pp(X=0,Y=1)=\frac12-\frac12=0.
\]
所以联合分布律为
\[
\begin{array}{c|cc}
 &Y=0&Y=1\\\hline
X=-1&\dfrac14&0\\
X=0&0&\dfrac12\\
X=1&\dfrac14&0
\end{array}
\]

\textbf{（2）求协方差：}

\[
\begin{aligned}
\E(X)&=(-1)\times\frac14+0\times\frac12+1\times\frac14=0,\\
\E(Y)&=0\times\frac12+1\times\frac12=\frac12.
\end{aligned}
\]
由于 $XY=0$ 几乎处处成立，故
\[
\E(XY)=0.
\]
于是
\[
\Cov(X,Y)=\E(XY)-\E(X)\E(Y)=0-0\times\frac12=0.
\]
}

\subsection*{4. 解答题}
\Q{1. 设二维随机变量 $(X,Y)$ 在平面区域 $G=\{(x,y)\mid0\le x\le2,\ 0\le y\le2\}$ 上服从均匀分布。求：（1）$(X,Y)$ 的联合概率密度；（2）$\Pp(X<1)$；（3）$\Pp(Y<1\mid X<1)$。}
\Ana{区域面积 $S_G=4$，故
\[
f(x,y)=\begin{cases}\dfrac14,&(x,y)\in G,\\0,&\text{其他}。\end{cases}
\]
\[
\Pp(X<1)=\frac{1\times2}{4}=\frac12.
\]
\[
\Pp(Y<1\mid X<1)=\frac{1\times1/4}{1/2}=\frac12.
\]}

\Q{2. 设总体 $X$ 的概率密度为
\[
f(x;\theta)=\begin{cases}\theta e^{-\theta x},&x>0,\\0,&x\le0,\end{cases}\qquad\theta>0,
\]
其中 $\theta$ 未知，$X_1,\ldots,X_n$ 是总体 $X$ 的样本。用极大似然估计法估计 $\theta$。}
\Ana{
\[
L(\theta)=\theta^n\exp\left(-\theta\sum_{i=1}^nX_i\right),
\]
\[
\ell(\theta)=n\ln\theta-\theta\sum_{i=1}^nX_i,
\]
\[
\ell'(\theta)=\frac n\theta-\sum_{i=1}^nX_i=0.
\]
\[
\hat\theta=\frac{n}{\sum_{i=1}^nX_i}=\frac1{\overline X}.
\]}

\subsection*{5. 应用题}
\Q{1. 小华从福州到广州出差，乘坐汽车、高铁、飞机的概率分别为 $0.2,0.3,0.5$；乘坐汽车、高铁、飞机时迟到的概率分别为 $0.2,0.1,0.04$。求：（1）小华迟到的概率；（2）已知小华迟到，求其乘坐汽车的概率。}
\Ana{设 $A_1,A_2,A_3$ 分别表示汽车、高铁、飞机，$B$ 表示迟到。
\[
\Pp(B)=0.2\times0.2+0.3\times0.1+0.5\times0.04=0.09.
\]
\[
\Pp(A_1\mid B)=\frac{0.2\times0.2}{0.09}=\frac49.
\]}

\Q{2. 人的智力商数（IQ）$X$ 服从正态分布 $N(100,15^2)$。求：（1）平均 IQ；（2）IQ 超过 $130$ 的概率；（3）IQ 在 $85$ 分至 $115$ 分的概率。已知 $\Phi(1)=0.8413$，$\Phi(2)=0.9772$。}
\Ana{
\[
\E(X)=100.
\]
\[
\Pp(X>130)=1-\Phi\left(\frac{130-100}{15}\right)=1-\Phi(2)=0.0228.
\]
\[
\Pp(85\le X\le115)=\Phi(1)-\Phi(-1)=2\Phi(1)-1=0.6826.
\]}

\subsection*{6. 证明题}
\Q{已知随机变量 $X\sim U(0,6)$。（1）求随机变量 $X$ 的概率密度；（2）证明 $\E[X(X-5)]=-3$。}
\Ana{
\[
f(x)=\begin{cases}\dfrac16,&0<x<6,\\0,&\text{其他}。\end{cases}
\]
由均匀分布性质，$\E(X)=3$，$\D(X)=3$，所以
\[
\E(X^2)=\D(X)+[\E(X)]^2=3+9=12.
\]
\[
\E[X(X-5)]=\E(X^2)-5\E(X)=12-15=-3.
\]}

\end{document}
